\documentclass[a4paper,11pt]{article}
\usepackage[czech]{babel}
\usepackage[T1]{fontenc}
\usepackage[utf8]{inputenc}
\usepackage{lmodern}
\usepackage{amsmath,amssymb}
\usepackage{geometry}
\usepackage{enumitem}
\geometry{a4paper,left=16mm,right=16mm,top=12mm,bottom=12mm}
\pagestyle{empty}
\setlength{\parindent}{0pt}
\setlength{\parskip}{2pt}
\setlist[enumerate]{leftmargin=7mm,itemsep=0.8mm,topsep=0.8mm}
\newcommand{\Title}[1]{%
\begin{center}{\large\bfseries ČTVRTLETNÍ TEST -- EXPONENCIÁLNÍ FUNKCE A ROVNICE}\\[1mm]
{\bfseries Řešení pro učitele -- varianta #1}\end{center}\hrule\vspace{3mm}}
\begin{document}
\Title{A}
\begin{enumerate}
\item \textbf{Základ exponenciální funkce.}
\begin{enumerate}[label=\alph*),leftmargin=6mm,itemsep=0.4mm,topsep=0.4mm]
\item $1^x=1$ pro všechna $x\in\mathbb R$, vznikne \textbf{konstantní funkce}.
\item Např. $(-4)^{1/2}=\sqrt{-4}\notin\mathbb R$; záporný základ nepřipouští všechny reálné exponenty.
\end{enumerate}
\item \textbf{Monotónnost a vlastnosti grafu.}
\begin{enumerate}[label=\alph*),leftmargin=6mm,itemsep=0.4mm,topsep=0.4mm]
\item Pro $a>1$ je $a^{-x}$ \textbf{klesající}, pro $0<a<1$ \textbf{rostoucí}; změna $x\mapsto -x$ obrací monotónnost.
\item Asymptota $\boxed{y=4}$, obor hodnot $\boxed{H(g)=(4;\infty)}$.
\end{enumerate}
\item \textbf{Rovnice a substituce.}
\begin{enumerate}[label=\alph*),leftmargin=6mm,itemsep=0.4mm,topsep=0.4mm]
\item Protože $5^{2-x}>0$ pro všechna $x$, je $\boxed{K=\varnothing}$.
\item Vyloučíme $t=-\frac12<0$, neboť $2^x>0$. Ze $2^x=1$ plyne $x=0$; $\boxed{K=\{0\}}$.
\end{enumerate}

\item Graf funkce $g(x)=2^{x-1}-2$ je rostoucí. Vodorovná asymptota je $\boxed{y=-2}$,
$\boxed{D(g)=\mathbb R}$ a $\boxed{H(g)=(-2;\infty)}$. Při náčrtu stačí vyznačit
asymptotu a např. body $[1;-1]$, $[2;0]$, $[3;2]$.
\item Položme $t=2^x>0$, pak $2^{2x+1}=2t^2$:
\[2t^2-9t+4=0,\quad D=81-32=49,\quad t_{1,2}=\frac{9\pm7}{4}.\]
$t_1=4$, $t_2=\frac12$; zpětně $2^x=4\Rightarrow x=2$,
$2^x=\frac12\Rightarrow x=-1$; $\boxed{K=\{-1;2\}}$.
\item Substituce $u=2^x>0$, $v=3^y>0$:
\[\begin{cases}u+5v=53,\\7u-v=47\end{cases}\quad\Longrightarrow\quad
\begin{cases}u=8,\\v=9.\end{cases}\]
Tedy $2^x=8\Rightarrow x=3$ a $3^y=9\Rightarrow y=2$;
$\boxed{K=\{[3;2]\}}$.
\end{enumerate}


\newpage
\Title{B}
\begin{enumerate}
\item \textbf{Základ exponenciální funkce.}
\begin{enumerate}[label=\alph*),leftmargin=6mm,itemsep=0.4mm,topsep=0.4mm]
\item $1^{x-2}-3=1-3=-2$, tedy \textbf{konstantní funkce}.
\item Pro $a>1$ je $a^x$ ryze rostoucí, pro $0<a<1$ ryze klesající; v obou případech je \textbf{prostá}.
\end{enumerate}
\item \textbf{Monotónnost a vlastnosti grafu.}
\begin{enumerate}[label=\alph*),leftmargin=6mm,itemsep=0.4mm,topsep=0.4mm]
\item Pro $a>1$ funkce $a^{-2x}$ \textbf{klesá}, pro $0<a<1$ \textbf{roste}; exponent $-2x$ mění směr.
\item Asymptota $\boxed{y=-5}$ a $\boxed{H(g)=(-5;\infty)}$.
\end{enumerate}
\item \textbf{Rovnice a substituce.}
\begin{enumerate}[label=\alph*),leftmargin=6mm,itemsep=0.4mm,topsep=0.4mm]
\item Jelikož $7^{x+1}>0$ pro každé $x\in\mathbb R$, nemůže se rovnat zápornému číslu $-\frac17$. Proto $\boxed{K=\varnothing}$. Pozor: $7^{-1}=\frac17$, nikoli $-\frac17$.
\item $4^x>0$, tedy $t=-2$ vyloučíme. $4^x=\frac14\Rightarrow x=-1$, $\boxed{K=\{-1\}}$.
\end{enumerate}

\item Graf $g(x)=\left(\frac13\right)^{x+1}-3$ je klesající. Vodorovná asymptota
$\boxed{y=-3}$, $\boxed{D(g)=\mathbb R}$, $\boxed{H(g)=(-3;\infty)}$.
K náčrtu stačí asymptota a např. body $[-2;0]$, $[-1;-2]$, $[0;-\frac83]$.
\item Nechť $t=3^x>0$, pak $3^{2x+1}=3t^2$:
\[3t^2-10t+3=0,\quad D=100-36=64,\quad t_{1,2}=\frac{10\pm8}{6}.\]
$t_1=3$, $t_2=\frac13$, tedy $x=1$ nebo $x=-1$;
$\boxed{K=\{-1;1\}}$.
\item Substituce $u=3^x>0$, $v=4^y>0$:
\[\begin{cases}2u+3v=24,\\3u-10v=7\end{cases}\quad\Longrightarrow\quad
\begin{cases}u=9,\\v=2.\end{cases}\]
Tedy $3^x=9\Rightarrow x=2$, $4^y=2\Rightarrow y=\frac12$;
$\boxed{K=\{[2;\frac12]\}}$.
\end{enumerate}


\end{document}
